\section{问题重述与分析}

光伏发电具有日内波动，居民负载与光伏出力的高峰也未必重合。并网微网可利用储能转移电量，并根据外网电价确定购电时段。其调控任务是在满足负载和储能物理约束的条件下，协调提前购电、储能运行与事后补购，降低实际支付费用。储能既可用于峰谷价差套利，也可用于吸收预测误差，两种用途竞争同一有限容量。

问题一给定单日电价、负载和光伏预测，要求每天零时形成购电计划，并使日首、日末储电量相同。该问题构成后续模型的物理基础：以时段购电量和储能充放电量为控制变量，以储电量为跨时段状态。日前储能调度及其线性建模可参考文献\cite{kong2016,silva2020}，本题的设备参数和约束数值均取自赛题。

问题二中负载和光伏实际值在决策时尚未知晓。计划买少会触发五倍电价的紧急购电，买多又不能获得退款，因此不能仅对预测均值重复求解问题一。本文用历史预测残差构造场景，将计划决策与紧急补救分层，并对高额紧急费用加入风险度量。两阶段随机调度与风险约束研究\cite{abunima2022,marino2018,grover2018}提供了这一结构的依据，但本题不引入文献中的柴油机、风电机组、可中断负载和售电收入。

问题三在每日0:00、6:00、12:00、18:00获得新光伏预报，可据此修订未执行的购电量。预测更新与调整费用必须同时考虑：预报更准确，并不自动意味着支付调整费用后仍然更省钱。本文把问题二的优化器置于定时滚动框架内，用实际储电量连接前后两次决策\cite{palma2013,silvente2015,li2016,jani2022}。问题四进一步将电价视为未知的未来交易价格，在保留物理约束与结算方式的基础上加入价格预测及联合残差场景\cite{li2016,li2020}。

四问采用共同的“预测、计划、执行、反馈”结构，递进增加不确定性和信息更新。模型的针对性设计是利用近期预测误差同时调节紧急费用风险权重与储能参考备用，并将优化参考轨迹和实际可执行控制分开。该设计是否节省费用，应以附件数据上的实际结算与对照计算判断，不能仅凭模型形式认定其优越性。

\section{模型假设与符号说明}

\subsection{基本假设}

\begin{enumerate}
  \item 将输入功率样本解释为对应10分钟时段的平均功率，以$\Delta t=1/6$小时换算电量。忽略微网内部线路损耗和电池自放电，电池充电、放电效率分别为$\eta_c=\eta_d=0.9$，故完整往返效率为0.81。
  \item 不向外网售电；超过负载与可充电能力的富余电量允许免费弃置。计划购电按承诺电量付费，未使用电量不退款；不额外引入题目未给出的寿命折损、碳排放或设备投资费用。
  \item 储能容量及功率限制按附录执行。问题一采用$e_0=e_{144}=6000$ kWh；问题二至四只在2025年1月1日零时初始化为6000 kWh，以后传递前一日实际末态，不每日重置。
  \item 外网可在当前时段补足剩余缺口，题目未给定紧急购电容量上限，故不另设此上限。由此实际负载能够得到满足，需优化的是补购费用与紧急购电发生情况，而不是宣称提高某一失供电概率指标。
  \item 在10分钟执行尺度上，可观测当前时段负载、光伏与储电量并即时平衡功率。未来实际值不可提前使用；历史价格、负载以及当前预报发布前已经结束时段的测量值均可用于决策。
  \item 附件4记录实际交易电价，不将当天全部价格提前告知日前优化器。预测和场景用于决策，实际价格用于结算，二者严格分离。
\end{enumerate}

\subsection{主要符号}

下文所有购电、充放电与储电变量均以kWh为单位，只有原始输入采用kW。为避免混淆，以$p$表示电价，以$c$表示充电量。

\begin{table}[htbp]
\centering
\caption{主要符号及含义}
\label{tab:symbols}
\small
\begin{tabular}{p{0.18\textwidth}p{0.62\textwidth}p{0.11\textwidth}}
\toprule
符号 & 含义 & 单位\\
\midrule
$D,t,s$ & 日期、日内时段、场景索引，$t=0,\ldots,143$ & ---\\
$L_{D,t},R_{D,t}$ & 实际负载功率与光伏发电功率 & kW\\
$n_{D,t}$ & 净负荷电量，$(L_{D,t}-R_{D,t})\Delta t$ & kWh\\
$p_{D,t}$ & 当前时段实际交易电价 & 元/kWh\\
$g_{D,t},a_{D,t}$ & 零时计划购电量、最终执行的调整购电量 & kWh\\
$c_{D,t},d_{D,t}$ & 电网侧充电输入量、放电输出量 & kWh\\
$e_{D,t}$ & 时段$t$开始时的储电量 & kWh\\
$h_{D,t},w_{D,t}$ & 紧急购电量、富余弃置电量 & kWh\\
$\pi_s,Q_s$ & 场景概率、场景整段紧急购电费用 & ---，元\\
$\rho_t,\lambda$ & 参考储能备用量、风险权重 & kWh，---\\
$E_{\min},E_{\max}$ & 储电量下限1200、上限10800 & kWh\\
$U$ & 单时段充放电量上限，$5000/6$ & kWh\\
\bottomrule
\end{tabular}
\end{table}

\section{数据处理与因果预测}

\subsection{时间尺度与结果表头约定}

附件2和附件4包含2025年365天、每天144个时刻的数据，附件3包含每天四次发布的24点整点光伏预报。对日期顺序、样本数、有限值和非负性逐项检查，原始负载、光伏、电价及预报均未发现缺失或负值；保留观测极值，不用全年统计量事后删除异常值。原始文件只读保存，并记录文件哈希以支持追溯。

源数据中的00:10作为右端点，对应物理区间00:00--00:10；最后一个标记$0{:}00{+}1$对应当日23:50--24:00。于是
\begin{equation}
 I_t=[t/6,(t+1)/6),\qquad
 n_{D,t}=(L_{D,t}-R_{D,t})/6.
 \label{eq:time_energy}
\end{equation}
最后一个区间的右端点24:00归入前一日结算。全文指定时段表均按上述自然日物理区间取数，四小时充放电量由24个时段求和。

附件5原模板的首列区间显示为0:10--0:20，末列显示为$0{:}00$--$0{:}10{+}1$，与上述物理区间存在10分钟显示偏移。按用户指定保留原表头及其格式，结果按源数据144列的顺序逐列一一填入：第一个数据列保存$I_0$结果，最后一列保存$I_{143}$结果。该规则是列序映射，不按模板区间字面再平移一次数据；工作簿显示标签不改变优化器的自然日边界。

\subsection{训练期划分与预测模型选择}

2025年1月用于建立预测、积累历史误差并推进储电状态；2月1日至12月31日的334天作为结果填报期。候选预测器包括前一日、3/7/14/28日移动平均、前一周、最近2/4个同星期日均值、同星期均值与近7日均值的加权组合，以及带近期水平修正的两种预测，共12种。对每个候选器，预测日期$D$时只能读取$D$以前的实际数据。

模型类型依据1月15日至31日逐日前推预测的均方根误差选择，2月以后固定类型，仅更新其可用历史窗口。负载和电价选中最近至多4个同星期日均值，历史光伏预测选中近7日均值。以$X$代表负载或电价，其预测形式为
\begin{equation}
 \widehat X_{D,t}=\frac{1}{|\mathcal W_D|}\sum_{j\in\mathcal W_D}X_{j,t},
 \quad \mathcal W_D=\{D-7,D-14,D-21,D-28\}\cap\{1,\ldots,D-1\},
 \label{eq:weekday_forecast}
\end{equation}
光伏则平均此前至多7天同一时段出力。历史不足时使用已有样本；首日用附件1作为固定先验。1月前15天使用近7日均值，此后仅按已经发生的1月预测误差在线选型，不把1月底才确定的模型选择结果反用于月初决策。

均方根误差用于选型，平均绝对误差用于补充描述，两者分别为
\begin{equation}
 \mathrm{RMSE}=\sqrt{\frac1N\sum_{i=1}^{N}(X_i-\widehat X_i)^2},
 \qquad \mathrm{MAE}=\frac1N\sum_{i=1}^{N}|X_i-\widehat X_i|.
\end{equation}
全年误差统计只评价所选方法，不参与重新选型。储能在1月也逐时执行并向后传递，因此2月1日初态是训练期运行的实际结果。

\subsection{发布预报插值与日内修正}

问题二不使用附件3，光伏完全依据历史预测；问题三及问题四对应的滚动策略使用当前已发布预报。将发布时刻记为$r\in\{0,6,12,18\}$小时，24个预报点解释为$r+1$至$r+24$小时。在发布时刻已观测光伏值与这些预报点之间作线性插值，得到10分钟右端点的预测；零时锚点取前一日24:00观测，首日零时取0。仅保留发布以后且不超过本日24:00的部分，跨日预报在当前实现中不进入优化。

日内负载预测以零时预测为基线，用此前2小时已结束的12个时段计算平均偏差$b_r$。设$k=1,2,\ldots$为发布后的预测步数，采用
\begin{equation}
 \widehat X_{r,k}=\max\!\left\{0,\widehat X^{(0)}_{r,k}
 +\operatorname{clip}\!\left(b_r\exp(-k/36),-0.3\widehat X^{(0)}_{r,k},0.3\widehat X^{(0)}_{r,k}\right)\right\}.
 \label{eq:issue_correction}
\end{equation}
其中$\operatorname{clip}$表示限制到给定上下界。6小时衰减尺度和30\%修正幅度预先固定；波动电价按同一规则修正，并设置正的数值下界。用于计算$b_r$的都是发布前已观测值。光伏直接采用本次发布预报，不使用事后实际曲线替代。
